\chapter{Introduction}
\label{ch:intro}

\section{What BAG is for}

A system on chip connects its processors, DMA engines and peripherals through bus
fabrics: crossbars, shared buses, bridges between protocols, and small adapters that
change a data width or cross a clock domain. Writing that interconnect by hand takes
time, and every change to the address map or the list of ports means more hand edits.

BAG, the Bus Architecture Generator, writes the interconnect for you. You describe it in
a \term{spec workbook}: an Excel file with one sheet per topic (masters, slaves, fabrics,
bridges, the memory map, and so on). BAG reads the workbook, checks it, and writes:

\begin{itemize}
\item the SystemVerilog of every fabric, bridge and adapter, inside a top-level
  \term{wrapper} module that your SoC instantiates;
\item a parameter package with the address map, a C header, and documents (memory map,
  port map, tie-offs, reset sequence);
\item cocotb tests that drive traffic through the generated design and check every
  route, response and arbitration rule;
\item SymbiYosys formal jobs that prove address decode and check the bus protocol on
  every port.
\end{itemize}

BAG supports four AMBA protocols: AXI4, AXI4-Lite, AHB-Lite and APB (APB3 and APB4).
The AXI fabrics are built on the open-source PULP \code{axi} library; the AHB-Lite and
APB fabrics, the bridges between protocol families and some adapters are written for
BAG.

\section{How BAG works}

BAG works in steps. Each step has its own command, and \cmd{bag flow} runs them in turn
(chapter~\ref{ch:cli}).

\begin{enumerate}
\item \textbf{Read and check.} \cmd{bag check} reads the workbook and runs every check
  (chapter~\ref{ch:checks}). An error stops here: no file is written.
\item \textbf{Resolve.} BAG derives everything you do not type: which fabric each port
  is on, the port counts, ID widths, the adapters to insert, the address decode tables.
  The result is the \term{IR} (intermediate representation), a YAML file that holds the
  whole design (chapter~\ref{ch:ir}). \cmd{bag ir} writes it.
\item \textbf{Generate.} \cmd{bag generate} writes the RTL, the collateral, the tests and
  the formal jobs into an \term{output directory} (chapter~\ref{ch:outputs}).
\item \textbf{Lint and compile.} \cmd{bag lint} runs slang, Verilator and Verible on the
  generated RTL; the design must have no finding that is not waived. A compile with
  Verilator follows.
\item \textbf{Simulate.} \cmd{bag sim} runs the generated cocotb tests with Verilator.
\item \textbf{Prove.} \cmd{bag formal} runs the generated SymbiYosys jobs.
\end{enumerate}

Steps 1 to 3 need only Python. Steps 4 to 6 need the EDA tools of
chapter~\ref{ch:install}.

\section{What you can rely on}

\begin{itemize}
\item \textbf{Same input, same output.} BAG writes the same bytes for the same workbook
  on every run and every machine. Generated files carry no time stamps. You can keep
  them under version control and review a change as a diff.
\item \textbf{The IR is the contract.} Every generated file comes from the IR. A
  workbook and an IR are both accepted as input; an IR is checked as strictly as a
  workbook.
\item \textbf{Every finding names its place.} A check message names the sheet, the row
  and a stable \term{rule ID}, for example \ruleid{TOPO-INDEX-OUT-OF-RANGE}.
\item \textbf{The generated files are yours.} They are not covered by the licence of
  BAG itself.
\end{itemize}

\section{What BAG does not do}

\begin{itemize}
\item It does not write your masters and slaves (processors, memories, peripherals).
  It writes the interconnect between them.
\item It does not run synthesis, place and route, or timing analysis. It writes timing
  constraints for clock-domain crossings, but no tool checks them.
\item It does not edit generated files for you. Change the workbook and generate again.
\item It does not test a fabric you supply yourself (a \code{bag_external} fabric): BAG
  does not know how it behaves.
\end{itemize}

Chapter~\ref{ch:limits} lists the known limits in detail.

\section{How to read this manual}

Chapter~\ref{ch:first} takes the sample design from an empty folder to a generated,
linted and simulated fabric. Read it first. Chapter~\ref{ch:workbook} is the reference
for every sheet and column of the workbook, chapter~\ref{ch:checks} for every check, and
chapter~\ref{ch:cli} for every command.

Every example in this manual is a file that a test runs, so the output you see is what
BAG prints. A line that starts with \code{$} is a command; the lines after it are what
the command prints. \code{<examples>} stands for the folder the example runs in, and
\code{<version>} for the version of BAG.
